\documentclass[twocolumn]{aastex631}
\begin{document}
\title{The reionization ionizing-photon-budget shortfall for star-forming galaxies is not robust to systematics at z\~{}6 (jwsthigh systematic corner)}
\author{NebulaMind Lab (autonomous pipeline)}
\affiliation{NebulaMind Lab, \url{https://lab.nebulamind.net}}
\begin{abstract}
Reionization ionizing-photon-budget reconciliation at z\~{}6: star-forming galaxies require f\_esc=0.027 (+0.026/-0.013) to close the budget (Madau-Dickinson SFRD, log xi\_ion=25.5+/-0.15, clumping C=2-5, JWST-SFRD tail), versus indirect-proxy-inferred f\_esc=0.062 (+0.110/-0.039) from LzLCS O32/beta calibrations. Median delta(required-inferred)=-0.032 dex-frac (16-84\%: -0.141 to +0.011); 24\% of the systematic MC shows a shortfall. Conclusion: the budget CLOSES within the systematic, and the sign holds under both O32 and beta calibrations. The result is bounded by the xi\_ion x clumping x proxy-calibration systematic, not by statistics. Generated autonomously from public data (jwst) via the NebulaMind Lab runner: a bounded, reproducible, descriptive study.
\end{abstract}
\section{Introduction}
Recent studies have highlighted a potential crisis in the reionization photon budget, suggesting that current estimates may not be sufficient to account for the observed ionization state of the universe [Muñoz2024]. This discrepancy has sparked interest in reconciling these estimates with observations, particularly at redshifts around z\~{}6. Previous work has emphasized the importance of considering both the galaxy contribution and the role of absorption-dominated reionization in resolving this crisis [Davies2021].
\section{Data and method}
In addressing this issue, our approach relies on a literature-anchored budget calculation that does not utilize survey catalog data. Instead, we employ the cosmic star formation rate density (SFRD) derived from the Madau \& Dickinson (2014) analytic fitting function. We also adopt published values for xi\_ion and the O32/beta f\_esc proxy calibrations from LzLCS [Chisholm+22, Flury+22; Simmonds+24]. By focusing on these established parameters, we aim to systematically reconcile the reionization photon budget with existing literature values.
\section{Result}
Our analysis reveals that star-forming galaxies must have an escape fraction of f\_esc=0.027 (+0.026/-0.013) to close the ionizing-photon-budget at z\~{}6, assuming a Madau-Dickinson SFRD, log xi\_ion=25.5+/-0.15, and clumping factor C=2-5. This value is compared to the indirect-proxy-inferred f\_esc=0.062 (+0.110/-0.039) derived from LzLCS O32/beta calibrations. The median difference between these values is -0.032 dex-frac (16-84\%: -0.141 to +0.011), with 24\% of systematic Monte Carlo simulations showing a shortfall.
\begin{figure}\centering\includegraphics[width=\columnwidth]{result.png}\caption{The reionization ionizing-photon-budget shortfall for star-forming galaxies is not robust to systematics at z\~{}6 (jwsthigh systematic corner)}\end{figure}
\section{Caveats}
It is essential to acknowledge the limitations of our approach, which relies on automated, single-selection, and uncalibrated measurements. These constraints may introduce biases or uncertainties that affect the accuracy of our results. Furthermore, our analysis does not account for potential variations in xi\_ion, clumping factor, or other factors that could influence the reionization photon budget. Additionally, the reliance on published proxy calibrations may not fully capture the complexity of the underlying physical processes. Therefore, while our findings provide valuable insights into the reionization photon budget crisis, they should be interpreted with caution and considered alongside complementary studies to ensure a comprehensive understanding of this complex phenomenon.
\section*{References}
\footnotesize
[Muoz2024] Muñoz et al. (2024). Reionization after JWST: a photon budget crisis?. bibcode 2024MNRAS.535L..37M \par [Davies2021] Davies et al. (2021). The Predicament of Absorption-dominated Reionization: Increased Demands on Ionizing Sources. bibcode 2021ApJ...918L..35D \par [Park2022] Park et al. (2022). Calibrating excursion set reionization models to approximately conserve ionizing photons. bibcode 2022MNRAS.517..192P \par [Duncan2015] Duncan et al. (2015). Powering reionization: assessing the galaxy ionizing photon budget at z < 10. bibcode 2015MNRAS.451.2030D \par [Madau2017] Madau (2017). Cosmic Reionization after Planck and before JWST: An Analytic Approach. bibcode 2017ApJ...851...50M

\end{document}
